\clearpage
\begin{summarybox}
\textbf{C091\quad 双曲线渐近线平行四边形面积定值——快记／快懂／快用}

\textbf{快记（看见就用）：}
\[
\boxed{H:\frac{x^2}{a^2}-\frac{y^2}{b^2}=1,\ a,b>0}
\qquad\Longrightarrow\qquad
\boxed{S=\frac{ab}{2}}.
\]
\textbf{前提：}从 $P\in H$ 分别作\textbf{两条渐近线}的平行线；
四边形由两条渐近线及这两条平行线围成。

\textbf{快懂（一个恒等式：差平方变乘积）：}
\[
Q=(u,ku),\ R=(v,-kv),\quad
P=Q+R,\qquad k=\frac ba.
\]
\[
\underbrace{x^2-\frac{y^2}{k^2}}_{a^2}
=\underbrace{(u+v)^2-(u-v)^2}_{4uv},\qquad
\underbrace{S}_{\text{面积}}=2k|uv|=\frac{ab}{2}.
\]
\textbf{只记主线：}双曲线把 $4uv$ \textbf{固定}，
平行四边形面积就是 $2k|uv|$，因此面积固定。

\textbf{初中起点（反比例函数的迁移）：}
$XY=K$ 使矩形面积为 $|K|$；
而沿两渐近线分解出 $uv=a^2/4$，
倾斜平行四边形的面积为 $2k|uv|$。
\textbf{不变的是乘积，不是两条边各自的长度。}

\textbf{升级结论（可直接解周长题）：}
\[
\boxed{|OQ|\,|OR|=\frac{c^2}{4}},\qquad
\boxed{L_{OQPR}\ge2c},\qquad c^2=a^2+b^2.
\]
周长等号当且仅当 $P=(\pm a,0)$（菱形）。
若\emph{固定 $c$、改变双曲线}，则
$\boxed{0<S\le c^2/4}$，最大值在 $a=b$ 时取得。



\textbf{快用（识别信号与题型）：}
\begin{itemize}
\item “\textbf{双曲线上的任意点＋分别平行两渐近线＋面积}”$\Rightarrow$立刻 $ab/2$；
\item 给 $S$ 求半轴参数，先写 $\boxed{ab=2S}$；
\item 若换成等值曲线 $x^2/a^2-y^2/b^2=t$，写 $\boxed{S=ab|t|/2}$；
\item 若还要求四边形顶点、边长，转用常规解析法求 $Q,R$。
\end{itemize}

\textbf{证明方法选择：}
\textbf{最快}——距离乘积除以渐近线夹角正弦，代入曲线方程；
\textbf{最低门槛}——设渐近线上的边向量 $(u,ku),(v,-kv)$，用平方差公式；
\textbf{常规}——解 $Q,R$ 两交点坐标并计算行列式；
\textbf{升级}——把 $P$ 写成两渐近线方向向量的线性组合，利用行列式的面积缩放。

\textbf{证明路线补全：}
\textbf{最低门槛}真正做到不使用行列式：用
竖直剪切 $(X,Y)\mapsto(X,Y-kX)$ 保面积，
把底、高变成 $|u|$ 与 $2k|v|$，直接得到 $S=2k|uv|$。
\textbf{更一般：}若独立一次式 $L_1,L_2$ 满足
$L_1L_2=K$，且系数行列式为 $D\ne0$，则 $S=|K|/|D|$。



\textbf{防错三句：} \textbf{$P,O$ 是对顶点；面积有 $1/2$；推广时取绝对值。}
\textbf{额外防错：}渐近线上构造退化，面积为零；
$S=ab/2$ 的逆命题需检查两条曲线
$x^2/a^2-y^2/b^2=\pm1$；$c$ 表示半焦距，焦距为 $2c$。


\textbf{可迁移思想：}当曲线方程是\emph{两个线性式的乘积等于常数}，
由相应两组平行线围成的平行四边形往往有面积不变量。
\end{summarybox}
